% 埃尔德什·帕尔（Paul Erdős）（综述）
% license CCBYSA3
% type Wiki

本文根据 CC-BY-SA 协议转载翻译自维基百科\href{https://en.wikipedia.org/wiki/Paul_Erd\%C5\%91s}{相关文章}。

\begin{figure}[ht]
\centering
\includegraphics[width=6cm]{./figures/04eef634120f8175.png}
\caption{} \label{fig_ARPR_1}
\end{figure}
保罗·埃尔德什（英语：Paul Erdős，匈牙利语：Erdős Pál [ˈɛrdøːʃ ˈpaːl]；1913 年 3 月 26 日—1996 年 9 月 20 日）是一位匈牙利数学家。他是 20 世纪最具产出力的数学家之一，也是提出最多数学猜想的人之一。\(^\text{[2][3]}\)埃尔德什在离散数学、图论、数论、数学分析、逼近论、集合论以及概率论等领域追求并提出了大量问题。\(^\text{[4]}\)他的大部分研究集中在离散数学领域，解决了其中许多长期未解的难题。他大力倡导并作出重要贡献的拉姆齐理论，研究的是在何种条件下秩序必然出现。总体而言，埃尔德什的工作更倾向于解决已存在的开放问题，而非开辟或探索全新的数学领域。他一生共发表约 1500 篇数学论文，这一数字至今无人能及。\(^\text{[5]}\

他以“社会化数学”的实践方式以及古怪的生活方式闻名；《时代》杂志称他为“怪人中的怪人”。\(^\text{[6]}\他坚信数学是一项社会性活动，因此过着漂泊不定的生活，唯一目的就是与其他数学家共同撰写论文。即使在晚年，他仍将清醒的每个时刻都献给数学研究；1996 年，他在华沙参加数学会议时去世。\(^\text{[7]}\

埃尔德什与合作者们的大量共同成果，促成了“埃尔德什数”的诞生——它表示某位数学家与埃尔德什之间的最短合著路径中的“步数”。
\subsection{生平}
保罗·埃尔德什于 1913 年 3 月 26 日出生在奥匈帝国的布达佩斯，\(^\text{[8]}\是安娜（婚前姓威廉，Wilhelm）与拉约什·埃尔德什（原姓英兰德，Engländer）[9][10] 的独子。他的两个姐姐在他出生前几天因猩红热去世，年仅三岁和五岁。\(^\text{[11]}\他的父母都是犹太人，任职于高中教授数学。埃尔德什从小便对数学表现出极大的兴趣。由于父亲在 1914 年至 1920 年间作为奥匈帝国的战俘被囚禁于西伯利亚，\(^\text{[10]}\他幼年时期主要由一位德国女家庭教师照顾，\(^\text{[12]}\母亲不得不长时间工作以维持生计。父亲在战俘营中自学了英语，但许多单词的发音不正确。后来父亲教儿子说英语时，保罗也继承了这种发音，并在其一生中都保持这种独特的口音。\(^\text{[13]}\

他通过阅读父母放在家中的数学书籍自学阅读。五岁时，只要告诉他一个人的年龄，他就能在脑中计算出这个人一生到目前为止大约活了多少秒。\(^\text{[12]}\由于两位姐姐早逝，他与母亲关系极为亲密，据说直到上大学前，他们一直同睡一张床。[14][15]

16 岁时，父亲向他介绍了两个将成为他终生热爱的主题——无穷级数与集合论。在高中时期，埃尔德什热衷于解答《KöMaL》（《中学数学与物理杂志》）每月刊登的数学问题，并成为其中的活跃解题者之一。\(^\text{[13]}\
